% ==============================================================================
% Chapter II — Related Work
% SABRE-Denoise: Complex-Valued 1D U-Net for NMR Signal Denoising
% Target: Journal of Magnetic Resonance
% ==============================================================================

\section{Related Work}
\label{sec:related_work}

SabreSDN sits at the intersection of four research threads: traditional NMR
signal processing, general-purpose deep-learning denoising, complex-valued
neural networks, and the emerging field of deep learning for NMR spectroscopy.
We review each in turn and identify the gaps that motivate our work.

\subsection{Traditional NMR Signal Processing for Denoising}
\label{sec:rw_traditional}

Signal-to-noise enhancement has been a central concern of NMR spectroscopy
since its inception. The classical toolkit operates primarily through
post-acquisition mathematical processing of the free induction decay (FID).

\textbf{Window functions (apodization).}
Multiplying the FID by a decaying window function---exponential, Gaussian,
sine-bell, or Lorentz--Gauss---before Fourier transformation is the most
widely used approach \cite{Ernst1987}. The exponential window (matched filter)
is optimal for sensitivity enhancement when the natural line-width is known:
it weights each time-domain point by $\exp(-\pi \cdot \text{LB} \cdot t)$,
maximizing the signal-to-noise ratio when LB matches the intrinsic
transverse relaxation rate $R_2^* = 1/(\pi \cdot \text{FWHM})$. However,
this optimality comes at a cost: the matched filter broadens the line-width
by exactly a factor of two (the ``sensitivity-resolution trade-off''), and
a mismatched LB parameter---inevitable when a spectrum contains peaks of
different line-widths---simultaneously degrades both sensitivity and
resolution. Lorentz--Gauss transformation partially decouples this trade-off
by combining Gaussian multiplication (resolution enhancement) with Lorentzian
broadening (sensitivity enhancement), but at the expense of introducing two
tunable parameters whose optimal values interact nonlinearly
\cite{Lindon1980}.

\textbf{Linear prediction and subspace methods.}
Linear prediction extrapolates the FID beyond the acquisition window using
an autoregressive model, effectively reducing truncation artifacts without
the sensitivity penalty of apodization \cite{Barkhuijsen1985, Koehl1994}.
Cadzow filtering \cite{Cadzow1988}, adapted to NMR by Hore and colleagues
\cite{Hore2013}, constructs a Hankel matrix from the FID and applies
singular value decomposition with rank truncation. This approach exploits
the fact that the deterministic NMR signal (a sum of exponentially damped
sinusoids) has a Hankel matrix of rank equal to the number of distinct
frequencies, whereas additive white noise fills all singular values uniformly.
Truncating to rank $r$ therefore suppresses noise while preserving signal
structure. Despite its elegance, Cadzow filtering has two practical
limitations: (1)~the signal rank $r$ is unknown \emph{a priori} and must
be estimated or tuned---overestimating $r$ retains noise, while
underestimating $r$ suppresses weak peaks; and (2)~SABRE-enhanced spectra
with widely varying peak intensities challenge the low-rank assumption,
as weak resonances may have singular values below the truncation threshold
even when $r$ is chosen correctly.

\textbf{Wavelet denoising.}
Wavelet-based methods \cite{Antonijevic2006, Gomez2016} threshold wavelet
coefficients in a transformed domain where signal and noise are more
separable. Their performance hinges on selecting an appropriate wavelet
basis and thresholding rule (hard vs.\ soft, universal vs.\ level-dependent).
While adaptive threshold selection algorithms exist (e.g., SURE, BayesShrink),
they operate on statistical criteria rather than on chemical knowledge of
peak positions and shapes.

\textbf{Common thread.}
All traditional methods share a fundamental architecture: they process
the NMR signal through a \emph{manually specified mathematical transform}
with one or more free parameters that must be tuned by an expert user for
each spectrum or dataset. They do not learn from data, and they encode
only generic signal priors (smoothness, low rank, sparsity) rather than
the specific spectral characteristics of a given chemical system under a
given hyperpolarization protocol. SabreSDN, by contrast, \emph{learns} the
denoising transformation from examples, and its architecture incorporates
NMR-specific inductive biases (complex-algebraic constraints, multi-scale
receptive fields, peak-region attention) that go beyond generic signal
priors.

\subsection{Deep Learning for Image and Signal Denoising}
\label{sec:rw_dl_denoising}

The past decade has seen deep convolutional neural networks (CNNs)
revolutionize image denoising. Several architectural innovations from this
literature directly inform our design.

\textbf{U-Net.}
Ronneberger et al.\ \cite{Ronneberger2015} introduced the U-Net for biomedical
image segmentation. Its symmetric encoder--decoder structure with skip
connections bridging corresponding resolution levels has proven remarkably
effective for dense prediction tasks where the output shares the spatial
structure of the input. The U-Net's inductive bias---local processing via
convolution, multi-scale representation via progressive down-sampling and
up-sampling, and fine-detail preservation via skip connections---maps
naturally onto the NMR denoising problem, where the output spectrum must
match the input in length and where both fine peak structure and broad
baseline trends must be simultaneously preserved. Our SDNBlock adopts the
U-Net backbone (Section~\ref{sec:architecture}) but makes it
\emph{complex-valued} throughout, adds \emph{domain-specific attention},
and replaces the bottleneck and skip connections with
\emph{NMR-tailored alternatives}.

\textbf{DnCNN and residual denoising.}
Zhang et al.\ \cite{Zhang2017} proposed DnCNN, which popularized residual
learning for denoising: rather than predicting the clean image directly,
the network predicts the noise residual, and the denoised output is obtained
by subtraction ($\hat{\mathbf{x}} = \mathbf{y} - \mathcal{R}(\mathbf{y})$).
Batch normalization and residual connections enable training deeper networks
(up to 17--20 layers) with stable gradient flow. Our global residual
connection (Equation~\ref{eq:overall_residual}) follows this principle.
However, DnCNN was designed for 2D natural images with real-valued pixels;
its direct 1D adaptation lacks several NMR-specific features: complex-valued
operations, multi-scale dilated convolutions tuned to NMR line-width ranges,
and attention mechanisms that identify spectral peaks.

\textbf{Beyond CNNs?}
Recent work has explored Transformer-based architectures \cite{Liang2021}
and diffusion probabilistic models \cite{Ho2020} for image denoising, often
achieving state-of-the-art results on benchmark datasets. For 1D NMR spectra
of length $N = 8192$, however, the $\mathcal{O}(N^2)$ complexity of
self-attention makes vanilla Transformers computationally prohibitive for
GPU training at batch sizes practical for denoising ($B \geq 128$).
Efficient attention variants (Linformer, Performer) reduce this complexity
but at the cost of introducing additional hyperparameters and losing the
convolutional inductive bias of translation equivariance, which is a natural
prior for spectra where a peak at one chemical shift should be processed
identically to a peak at another. Diffusion models, while powerful, require
iterative sampling (typically 50--1000 forward passes) for a single denoised
output, which is incompatible with the real-time or near-real-time throughput
expected in an NMR acquisition pipeline. For these reasons, we build on the
mature and efficient U-Net backbone, enhanced with NMR-specific components.

\subsection{Complex-Valued Neural Networks}
\label{sec:rw_complex}

The SabreSDN architecture depends critically on complex-valued convolutions.
We trace this capability to the deep complex networks literature and
distinguish our contribution from prior applications.

\textbf{Foundations.}
Trabelsi et al.\ \cite{Trabelsi2018} established the modern framework for
deep complex networks, providing complex formulations of convolution,
batch normalization, activation functions, and weight initialization.
The complex convolution rule
$\mathbf{W} * \mathbf{x} = (\mathbf{A} * \mathbf{x}_r - \mathbf{B} * \mathbf{x}_i) +
j(\mathbf{B} * \mathbf{x}_r + \mathbf{A} * \mathbf{x}_i)$
that we use in \texttt{CConv1d} (Section~\ref{sec:cconv}) is a direct
implementation of their formulation. Our contribution is not the complex
convolution rule itself, but rather its \emph{application to a domain where
the complex structure is physically meaningful}: in NMR spectroscopy, the
real and imaginary components are not arbitrary channels but are related
through the Hilbert transform (the Kramers--Kronig relations for causal
linear response). This means that the complex-algebraic weight-sharing
constraint in \texttt{CConv1d} aligns with a genuine physical prior, whereas
in the classification and speech tasks explored by Trabelsi et al., the
complex representation is synthetically constructed from real-valued inputs
and the constraint serves primarily as a parameter-efficiency regularizer.

\textbf{MRI reconstruction.}
Complex-valued U-Nets have been successfully applied to magnetic resonance
image reconstruction from undersampled $k$-space data
\cite{Virtue2017, Cole2021, Sriram2020}. These works operate in 2D (image
domain) and exploit the fact that MRI data are natively complex-valued.
Our scenario differs in two key respects: (1)~we operate on 1D frequency-domain
spectra rather than 2D $k$-space data, meaning our spatial axis represents
chemical shift rather than spatial position; and (2)~the absorption--dispersion
relationship in high-resolution NMR is a strict Hilbert transform pair for
Lorentzian lines, providing a stronger physical constraint than the general
complex structure of MRI images. To our knowledge, SabreSDN is the first
application of complex-valued deep learning to 1D NMR spectral denoising.

\textbf{Why not just use real convolutions?}
A skeptic might ask whether complex convolutions are necessary: an
unconstrained two-channel real convolution can, in principle, approximate
any complex-linear operation (the universal approximation property of neural
networks guarantees this). The practical answer has two parts. First,
\emph{sample efficiency}: by hard-coding the complex multiplication structure,
\texttt{CConv1d} reduces the number of free parameters in each convolutional
layer by a factor of 2 (each complex weight pair $[w_r, w_j]$ maps to four
real convolution operations with shared parameters, vs.\ four independent
real convolutions in the unconstrained case), and this reduction translates
to lower sample complexity---the network needs fewer training examples to
achieve the same generalization performance. This is particularly relevant
in NMR, where the number of distinct clean spectra available for training
may be limited (tens rather than thousands). Second, \emph{physical
interpretability}: the complex-algebraic structure ensures that the
absorption and dispersion channels are processed in a manner consistent
with the physics of the NMR signal, which constrains the hypothesis space
to physically plausible transformations and reduces the risk of learning
spurious correlations that do not generalize.

\subsection{Deep Learning for NMR Spectroscopy}
\label{sec:rw_nmr_dl}

The application of deep learning to NMR spectroscopy is a rapidly growing but
still nascent field. Existing work clusters around three tasks:
identification, quantification, and reconstruction---of which the first two
have received substantially more attention.

\textbf{Metabolite identification and peak picking.}
NMRNet \cite{Le2019} and related architectures use CNNs to identify
metabolites in complex biological mixtures from 1D $^1$H NMR spectra.
DeepPicker \cite{Klukowski2020} and DEEP Picker \cite{Li2021} apply
deep learning to the peak-picking problem, detecting resonance positions
in crowded spectra. These methods operate at the \emph{analysis layer}:
their input is a spectrum of sufficient quality that peaks are visually
(or at least statistically) discernible, and their output is a higher-level
chemical interpretation (compound identity, peak list).

\textbf{Signal enhancement and reconstruction.}
Closer to our work, several groups have explored deep learning for NMR signal
enhancement. Non-uniform sampling (NUS) reconstruction using deep learning
\cite{Hyberts2012, Qu2015, Hansen2019} learns a mapping from sparsely
sampled FIDs to fully sampled spectra, reducing acquisition time rather
than suppressing noise per se. Deep neural networks have been applied to
virtual decoupling \cite{Karunanithy2021} and to the removal of experimental
artifacts \cite{Cobas2024}. These works demonstrate that deep learning can
learn meaningful transformations of NMR data, but they address different
noise or artifact sources than the broadband thermal and instrumental noise
that SabreSDN targets.

\textbf{The gap.}
To our knowledge, no prior work proposes a \emph{dedicated denoising
architecture for SABRE-hyperpolarized NMR spectra} that operates directly
in the complex domain, incorporates multi-scale receptive fields tuned to
the range of NMR line-widths, and uses a composite loss function that
prioritizes the spectral regions relevant to quantitative analysis.
SabreSDN fills this gap at the \emph{signal recovery layer}---downstream
of acquisition but upstream of any chemical analysis. If successful, it
provides a drop-in preprocessing module that can improve the input quality
for all downstream tasks (peak picking, integration, metabolomics, and
quantification) without modifying those downstream pipelines.

\subsection{Attention Mechanisms in Deep Learning}
\label{sec:rw_attention}

Attention mechanisms enable neural networks to selectively emphasize
informative features while suppressing irrelevant ones. Their adoption
in computer vision has progressed from channel-wise recalibration to
sophisticated spatial and hybrid schemes.

\textbf{Channel attention.}
Hu et al.\ \cite{Hu2018} introduced the Squeeze-and-Excitation (SE) block,
which compresses global spatial information into a channel descriptor via
average pooling, then learns channel-wise scaling factors through a
bottleneck fully-connected network. This mechanism has been shown to
improve representation quality across a wide range of architectures at
negligible computational cost. Our Channel Attention module
(Section~\ref{sec:cbam}) is a direct 1D adaptation of the SE design.

\textbf{Spatial and spectral attention.}
Woo et al.\ \cite{Woo2018} proposed the Convolutional Block Attention Module
(CBAM), which augments channel attention with a spatial attention branch
that identifies informative spatial regions via a small-kernel convolution.
Our \texttt{CBAM1d} is a dimension-reduced adaptation: the original
$7 \times 7$ 2D spatial convolution becomes a $7 \times 1$ 1D convolution
along the spectral axis, transforming the semantic meaning from
``which image region is salient'' to ``which chemical shift is occupied by
a peak.'' This reinterpretation is not merely a dimensionality reduction;
it encodes the prior that \emph{peak position is the primary axis of
variation in a 1D spectrum}, in contrast to natural images where
information is distributed across two spatial dimensions.

\textbf{Why CBAM and not self-attention?}
Self-attention \cite{Vaswani2017} and its vision variants \cite{Dosovitskiy2021}
have achieved dominant performance in many domains, raising the natural
question of why we retain a convolutional attention module rather than
adopting a Transformer backbone. Our rationale is threefold:
(1)~\textbf{Computational tractability}: the $\mathcal{O}(N^2)$ complexity of
self-attention over $N = 8192$ spectral points requires
${\sim}67 \times 10^6$ pairwise interactions per attention head, per layer,
making training at batch sizes of 128--256 infeasible on a single laboratory
GPU.
(2)~\textbf{Inductive bias alignment}: NMR denoising is fundamentally a
\emph{local-to-local} mapping---the clean signal at chemical shift $\delta$
depends overwhelmingly on the noisy signal at $\delta$ and its immediate
neighborhood, not on distant spectral regions where unrelated peaks reside.
Convolutional operations with their built-in locality and translation
equivariance are a near-perfect match for this structure, whereas global
self-attention introduces unnecessary degrees of freedom that the network
must learn to ignore.
(3)~\textbf{Channel attention already provides global context}: the global
average pooling in the Channel Attention branch aggregates information from
the entire spectrum into a per-channel descriptor, giving the network a
mechanism to condition local processing on global spectral statistics
(e.g., overall noise level, presence or absence of a dominant solvent peak)
without the quadratic cost of explicit pairwise attention.

\subsection{Positioning of SabreSDN}
\label{sec:rw_positioning}

Figure~\ref{fig:positioning} schematically positions SabreSDN within the
NMR deep learning landscape. The horizontal axis spans from generic signal
processing to NMR-specific methodology; the vertical axis spans from
analysis-layer tasks (identification, classification) to signal-layer tasks
(denoising, reconstruction, artifact removal).

\begin{figure}[t]
    \centering
    % \includegraphics[width=0.8\columnwidth]{figures/positioning.pdf}
    \caption{[TODO] Schematic positioning of SabreSDN in the NMR deep learning
             landscape. Horizontal axis: generic $\to$ NMR-specific methodology.
             Vertical axis: signal-recovery $\to$ chemical-analysis tasks.
             SabreSDN occupies the NMR-specific $\times$ signal-recovery
             quadrant.}
    \label{fig:positioning}
\end{figure}

SabreSDN occupies the \textbf{NMR-specific $\times$ signal-recovery} quadrant:
it is built on a generic deep learning backbone (U-Net) but customized at
every level to the physics and analytical priorities of SABRE NMR
spectroscopy---complex convolutions for absorption--dispersion coupling,
multi-scale bottleneck for heterogeneous line-widths, peak-focused attention
for quantitative fidelity, and an asymmetric loss function that prioritizes
the spectral regions that chemists actually care about.

In summary, SabreSDN draws on established ideas from each of the four
research threads reviewed above, but their synthesis into a single
architecture designed \emph{from first principles} for 1D complex NMR
spectral denoising is, to our knowledge, novel.
